\section*{Chapter \ref{ch:s1:carry-across-asset-classes} --- Carry Across Asset Classes}

\begin{solution}{exo:s1:carry-across-asset-classes:1}
$\ln(100/98) \times 252/21 = 0.242$: 24.2\% a year, a steeply backwardated curve.
\end{solution}

\begin{solution}{exo:s1:carry-across-asset-classes:2}
$5\% - 1\% = 4\%$ a year for the currency; $2\% - 4.5\% = -2.5\%$ a year for the index future: a long equity future pays the financing rate and
receives the dividends.
\end{solution}

\begin{solution}{exo:s1:carry-across-asset-classes:3}
Its spot price has an annual cycle that the curve anticipates: the one-month slope measures mostly where the season is going (winter prices above
summer ones), not the carry that would be earned if the whole curve stayed put. On the synthetic seasonal commodities the one-month slope has a
correlation of 0.05 with the true carry; a twelve-month slope compares the same season and cancels the cycle.
\end{solution}

\begin{solution}{exo:s1:carry-across-asset-classes:4}
A hundred-day standard deviation is $10\% \times \sqrt{100/252} = 6.3\%$, so the loss is $26.3/6.3 = 4.2$ standard deviations. It expected $0.95 \times
10\% \times 100/252 = 3.8\%$.
\end{solution}

\begin{solution}{exo:s1:carry-across-asset-classes:5}
Diversification relies on the classes' low correlation, which holds on average; in the crash every class lost at once, so the diversified book, scaled
to the same 10\% volatility on its low average correlation, holds more gross risk than any one class and loses more.
\end{solution}

\begin{solution}{exo:s1:carry-across-asset-classes:6}
Expected excess return is carry plus expected price appreciation. Carry is known in advance; the question is how much of it prices take back. If none,
the carry book earns its whole carry spread; if prices moved to cancel it, as uncovered interest parity says, it earns nothing. In the synthetic
universe the Sharpe ratios roughly double or triple between $\lambda = 0.5$ and $\lambda = 1$.
\end{solution}

\begin{solution}{exo:s1:carry-across-asset-classes:7}
The commodity book falls from 0.59 to 0.34 and the diversified book from 0.95 to 0.82: for the three seasonal commodities the one-month slope ranks the
season instead of the carry, so a third of the commodity book trades noise.
\end{solution}

\begin{solution}{exo:s1:carry-across-asset-classes:8}
Ten calm years say little about a strategy whose losses come in rare crashes: the synthetic currency book at 10\% volatility lost 25--29\% in each
crash. Measure the skewness and the exposure to global risk, stress the book with past unwinds, and size it for the crash, not the calm.
\end{solution}

\begin{solution}{pb:s1:carry-across-asset-classes:1}
\begin{enumerate}
\item The return if prices stay unchanged, $(S - F)/(F\tau)$; forward discount, dividend yield minus financing, yield plus roll-down, roll yield.
\item Expected excess return equals carry plus expected price appreciation.
\item Long high-carry and short low-carry assets of a class by rank; carry books of several classes at equal risk.
\item As the log slope between the first two contracts; for seasonal commodities the slope measures the season (correlation 0.05), so a twelve-month
slope is used.
\item Centred ranks with a gross of one, 2, 1, 2 and 4 basis points per unit traded, each class scaled in sample to 10\%; the diversified book equally
weighted and scaled the same way.
\item $\lambda = 0.5$: 0.17, 0.22, 0.03, 0.27 and 0.34; $\lambda = 1$: 0.28, 0.53, 0.49, 0.59 and 0.95.
\item Four books with low average correlation combine into one with a Sharpe ratio about twice the average class's.
\item Its crashes take back the carry it earns.
\item A sudden loss when high-carry assets fall together; negatively skewed moves between high- and low-rate currencies, from unwinding when funding
liquidity falls.
\item High-rate currencies load more on a global risk factor, so carry investors hold global risk in bad times.
\item Currency: 29.0\%, 24.9\% and 28.0\%; diversified: 12.9\%, 26.3\% and 13.1\% ($\lambda = 1$).
\item The unwind of currency carry trades in early August 2024 and the short-lived rise of the yen.
\item \textbf{Named result.} With all carry earned the diversified book's Sharpe ratio is 0.95 against 0.28--0.59 for single classes; in the second
crash it lost 26.3\% at 10\% volatility, more than any class alone (currencies 24.9\%, equities 12.4\%, commodities 11.3\%, bonds 4.1\%).
\item Their returns are uncorrelated on average and opposite in crashes: trend gains, carry loses.
\item An average of 1.2\% a year, backwardation on 51.7\% of days; 24.8\% a year at a Sharpe ratio of 0.61.
\item $+126\%$ on 14 January 1991 and $-575\%$ on 21 April 2020, the day after the May contract settled below zero.
\item Carry from prices known at the time, the contracts actually held, costs, and crash periods in the sample.
\item WTI curve carry.
\item Carry strategies of all classes do poorly together in global recessions.
\item For holding assets that lose in bad times: global risk and crash risk.
\end{enumerate}
\end{solution}

\begin{solution}{iq:s1:carry-across-asset-classes:1}
The interest differential (forward discount); the yield over the financing rate plus the roll-down; the roll yield, the slope of the curve from the spot
or nearby to the contract held.
\end{solution}

\begin{solution}{iq:s1:carry-across-asset-classes:2}
Uncovered interest parity says high-rate currencies depreciate by the rate difference on average, so the carry is taken back. Empirically they have
depreciated by less, often not at all on average, with occasional crashes; the difference is the carry premium.
\end{solution}

\begin{solution}{iq:s1:carry-across-asset-classes:3}
Recognise the start of an unwind: cut gross exposure to limits set in advance, starting with the most crowded pairs, rather than waiting for losses to
force it; avoid adding into the move; reassess when volatility settles.
\end{solution}

\begin{solution}{iq:s1:carry-across-asset-classes:4}
With stress tests on past unwinds applied to today's positions, the skewness of its returns in longer samples and other markets, its beta to global
risk factors, and the positioning of other carry traders.
\end{solution}

\begin{solution}{iq:s1:carry-across-asset-classes:5}
Forward points and interest rates, dividend forecasts and index futures prices, bond curves and cheapest-to-deliver data, commodity curves by contract,
roll calendars and volatility forecasts; stale quotes, missing contracts, holidays, roll errors and currency conversions can all corrupt a carry.
\end{solution}

\begin{solution}{iq:s1:carry-across-asset-classes:6}
$\ln F = \ln S - c\tau$ and $d\tau = -dt$, so $dF/F = dS/S + c\,dt$ (plus a second-order term that vanishes in expectation with $\ln$). With $S$ a
martingale, $E[dF/F] = c\,dt$. With $E[dS/S] = -(1 - \lambda)c\,dt$, $E[dF/F] = \lambda c\,dt$: the share $\lambda$ of the carry is earned.
\end{solution}
